\part{Cosmology}

\chapter{Friedmann-Robertson-Walker}

    Cosmology is the study of the Universe as a whole. Since the only non-negligible interaction for the study of the large-scale structure of the Universe is gravity, we can study the Cosmos by means of the general theory of relativity. In this chapter, after some reasonable assumptions of isotropy and homogeneity for the cosmic fluid, we will derive a suitable metric to mathematically describe the Universe, the so-called Friedmann-Robertson-Walker or FRW metric, and find the corresponding Einstein field equations, the Friedmann equations. Besides, we will delve into the solutions of these equations in the case of open, flat or close Universe.

\section{Principles}
    
    Our starting point is two principles upon which modern cosmology is based.
    \begin{princ}[Copernican]
        There is no preferred observer in the Universe.
    \end{princ}
    \begin{princ}[Cosmological]
        The Universe is homogeneous and isotropic.
    \end{princ}
    \noindent Actually, the former is implied by homogeneity of the latter: if there are no preferred observers, the Universe looks the same from any point, thus they can all be considered equivalent. Every assumption should be experimentally or observationally justified. Of course, isotropy could be tested by looking at the night sky, but what about homogeneity? Actually, even the night sky looks everything but isotropic, but we can restrict our attention to little sampled areas and checking the average distribution of matter in all directions. In addition, we have to take into account the fact that the image that we see at a given instant of time is in reality is the past light-cone. This implies that the Cosmological principle must be reformulated. 
    \begin{princ}[Cosmological II]
        There exists a time $t$ such that the Universe is homogeneous and isotropic on each time slice $\Sigma_t$ of it.
    \end{princ}
    \noindent Therefore, the Cosmological principle need to be considered as a working hypothesis and to be verified a posteriori.

\section{Symmetries}

    The metric has $10$ independent components, but it can be much simplified by studying the symmetries of the system. Isotropy entails the existence of three rotational space-like Killing vectors $K_a$ with $a=1,2,3$ such that they satisfy the Lie algebra $[K_a, K_b] = 0$; homogeneity entails the existence of three translational space-like Killing vectors $K_i$ with $i=1,2,3$ such that they satisfy the Lie algebra $[K_i, K_j] = \epsilon_{ijk} K_k$. We are in the so-called maximal symmetry and there are no more isometries in space. In fact, it can be proved that isotropy with respect to an arbitrary point is equivalent to homogeneity. Concerning time, there is no Killing vector associated to it. Accordingly, the metric depends only on a function of (coordinate) time $a=a(t)$, called the cosmic scale factor, to be determined by the Einstein field equations. Adopting as comoving coordinates $(t, r, \theta, \phi)$, the metric explicitly looks like 
    \begin{equation*}
        ds^2 = - dt^2 + a^2 (t) \Big ( \frac{dr^2}{1 - k r^2} + r^2 d\Omega^2 \Big ) ~.
    \end{equation*}
    This is known as the Friedmann-Robertson-Walker or FRW metric. Notice that there is a factor $k$, called the curvature scalar, that in principle could take any real value. However, it is always possible to rescale the $r$ coordinate in order to have $r = \tilde r / |k|^{1/2}$, so that 
    \begin{equation*}
    \begin{aligned}
        ds^2 & = - dt^2 + a^2 (t) \Big ( \frac{d r^2}{1 - k r^2} + r^2 d\Omega^2 \Big ) \\ & = - dt^2 + a^2 (t) \Big ( \frac{d \tilde r^2 / |k|}{1 - \tilde r^2 k / |k|} + \frac{\tilde r^2}{|k|} d\Omega^2 \Big ) \\ & = - dt^2 + \frac{a^2 (t)}{|k|} \Big ( \frac{d r^2}{1 - \tilde r^2 k / |k|} + r^2 d\Omega^2 \Big ) \\ & = - dt^2 + \tilde a^2 (t) \Big ( \frac{d r^2}{1 - \tilde r^2 \tilde k} + r^2 d\Omega^2 \Big ) ~,
    \end{aligned}
    \end{equation*}
    where we have rescaled $a$ by $\tilde a = a / |k|$ and $\tilde k = k / |k|$. In this way, the only values that $k$ can take are $0, \pm 1$. To understand the physical meaning of these values, let us make a change of variable of $r$ into $\psi$ such that 
    \begin{equation*}
        ds^2 = - dt^2 + a^2(t) (d\psi^2 + f_k(\psi)^2 d\Omega^2 ) ~, 
    \end{equation*}
    where 
    \begin{equation*}
        f_k(\psi) = \begin{cases}
            \psi & k=0 \\
            \sin \psi & k=1 \\
            \sinh \psi & k=-1 \\
        \end{cases} ~.
    \end{equation*}
    Depending on the sign of the curvature scalar $k$, we have 
    \begin{enumerate}
        \item if $k=0$, a flat universe with metric
            \begin{equation*}
                ds^2 = - dt^2 + a^2(t) (dr^2 + r^2 d\Omega^2 ) = - dt^2 + a^2(t) (d\psi^2 + \psi^2 d\Omega^2 ) 
            \end{equation*}
            and hypersurface at constant time $\Sigma_t$ to be flat;
        \item if $k=1$, an open universe with metric
            \begin{equation*}
                ds^2 = - dt^2 + a^2 (t) \Big ( \frac{dr^2}{1 - r^2} + r^2 d\Omega^2 \Big ) = - dt^2 + a^2(t) (d\psi^2 + \sin^2 \psi d\Omega^2 ) 
            \end{equation*}
            and hypersurface at constant time $\Sigma_t$ to be a $3$-dimensional sphere;
        \item if $k=-1$, a closed universe with metric
            \begin{equation*}
                ds^2 = - dt^2 + a^2 (t) \Big ( \frac{dr^2}{1 + r^2} + r^2 d\Omega^2 \Big ) = - dt^2 + a^2(t) (d\psi^2 + \sinh^2 \psi d\Omega^2 ) 
            \end{equation*}
            and hypersurface at constant time $\Sigma_t$ to be a $3$-dimensional hyperboloid.
    \end{enumerate}
    Now, we can express which is the meaning of the coordinates. The areal radius, i.e.~the radius such that the area is given by $A = 4 \pi r^2_a$, is given by 
    \begin{equation*}
        r_a = a(t) r ~,
    \end{equation*}
    whereas the proper distance is
    \begin{equation*}
        dR = \sqrt{g_{00}} dr = \frac{a(t)}{\sqrt{1 - k r^2}} dr ~.
    \end{equation*}
    Notice that both distances depend on time. $a$ gives the size of the Universe: for $k=1$, it is the radius of the Universe-sphere; for $k=0$, we can substitute the flat Universe with a torus of size $a$; for $k=-1$, even though such interpretation is not possible, it gives the size of the curvature of the Universe-hyperboloid. Therefore, as observations suggest, the Universe is expanding through $a(t)$.

\section{Computations}

    In matrix notation, the metric takes the form 
    \begin{equation*}
        \pyc{FRW.printmet()} ~.
    \end{equation*}
    The non-vanishing Christoffel symbols are 
    \begin{equation*}
        \pyc{FRW.printchri()} ~.
    \end{equation*}
    The non-vanishing Einstein tensor components are 
    \begin{equation*}
        \pyc{FRW.printeinup()} ~.
    \end{equation*}

\section{Cosmic fluid}

    Having identified the geometrical part of the system, i.e.~the metric, it is time to study the matter content. Our assumption here is that the Universe is a perfect fluid, whose energy-momentum tensor looks like 
    \begin{equation*}
        T_{\mu\nu} (t) = (p (t) + \rho(t) ) u_\mu u_\nu + p (t) g_{\mu\nu} ~, \quad T^{\mu}_{\phantom \mu \nu} (t) = \diag (- \rho(t), p(t), p(t), p(t)) ~,
    \end{equation*}
    where $u_\mu$ is the comoving $4$-velocity, i.e.~the fluid is at rest with it, $\rho(t)$ is the energy density and $p(t)$ is the pressure. The $0$-th component of the continuity equation for the energy-momentum tensor is the energy conservation 
    \begin{equation*}
        \dot \rho + 3 H (p + \rho) = 0 ~,
    \end{equation*}
    where $H (t) = \dot a(t) / a(t)$ is the Hubble constant. 
    \begin{proof}
        In fact, the $0$-th component is given by
        \begin{equation*}
        \begin{aligned}
            0 & = \nabla_\mu T^{\mu}_{\phantom \mu 0} = \partial_\mu T^{\mu}_{\phantom \mu 0} + \Gamma^\mu_{\nu \mu} T^{\nu}_{\phantom \nu 0} - \Gamma^\nu_{0 \mu} T^{\mu}_{\phantom \mu \nu} \\ & = \partial_0 T^{0}_{\phantom 0 0} + \Gamma^1_{0 1} T^{0}_{\phantom 0 0} + \Gamma^2_{0 2} T^{0}_{\phantom 0 0} + \Gamma^3_{0 3} T^{0}_{\phantom 0 0} - \Gamma^1_{0 1} T^{1}_{\phantom 11} - \Gamma^2_{0 2} T^{2}_{\phantom 2 2} - \Gamma^3_{0 3} T^{3}_{\phantom 33} ~.
        \end{aligned}
        \end{equation*}
        Now, we use the non-vanishing Christoffel symbols 
        \begin{equation*}
            \Gamma^1_{0 1} = \Gamma^2_{0 2} = \Gamma^3_{0 3} = \frac{\dot a}{a}  
        \end{equation*}
        to obtain
        \begin{equation*}
            0 = - \dot \rho - 3 \frac{\dot a}{a}  \rho - 3 \frac{\dot a}{a} p =  - \dot \rho - 3 H \rho - 3 H p = - \dot \rho - 3 H (\rho + p) ~.
        \end{equation*}
    \end{proof}
    \noindent In order to prove that it is indeed an energy conservation law, notice that we can define a comoving volume $V(t) \sim a^3(t)$ such that the total comoving energy is $E(t) = \rho(t) V(t)$. In fact
    \begin{equation*}
    \begin{aligned}
        \dv{E}{t} & = \dv{}{t} (\rho V) = \dv{}{t} (\rho a^3) = \dot \rho a^3 + \rho \dv{}{t} a^3 = - 3 \frac{\dot a}{a} (\rho + p) a^3 + 3 \rho \dot a a^2 \\ & = - 3 \frac{\dot a}{a} (\rho + p) a^3 + 3 \rho \frac{\dot a}{a} a^3 = - 3 \frac{\dot a}{a} p a^3 = - p \dv{}{t} a^3 = - p \dv{V}{t} ~,
    \end{aligned}
    \end{equation*}
    which is the thermodynamics relation $dE = - p dV$.

    In addition, we need to specify an equation of state, which relate pressure and energy density $p = p (\rho)$ for a barytropic fluid. The most simple relation we can have is linear 
    \begin{equation*}
        p = \omega \rho ~,
    \end{equation*}
    for some constant parameter $\omega$ to be determined. The solution of the energy conservation is 
    \begin{equation*}
        \rho \propto a^{- 3 (1 + \omega)} ~.
    \end{equation*}
    \begin{proof}
        In fact, 
        \begin{equation*}
            0 = \dot \rho + 3 H (\rho + p) = \dot \rho + 3 H \rho (1 + \omega) ~,
        \end{equation*}
        \begin{equation*}
            \frac{\dot \rho}{\rho} = - 3 (1 + \omega) \frac{\dot a}{a} ~, \quad 
            \frac{d \rho}{\rho} = - 3 (1 + \omega) \frac{d a}{a} ~,
        \end{equation*}
        \begin{equation*}
            \int \frac{d \rho}{\rho} = - 3 (1 + \omega) \int \frac{d a}{a} ~, \quad 
            \ln \rho \propto - 3 (1 + \omega) \ln a ~,
        \end{equation*}
        \begin{equation*}
            \ln \rho \propto \ln a^{- 3 (1 + \omega)} ~, \quad 
            \rho \propto a^{- 3 (1 + \omega)} ~.
        \end{equation*}
    \end{proof}
    The cosmic fluid can be decomposed into two part: dust, i.e.~non-relativistic pressure-less particles, and radiation, i.e.~highly-relativistic particles 
    \begin{enumerate}
        \item dust, i.e.~non-relativistic pressure-less particles, with $\omega = 0$ and $\rho \propto a^{-3}$;
        \item radiation, i.e.~non-relativistic pressure-less particles, with $\omega = 1/3$ and $\rho \propto a^{-4}$;
        \item dark energy, compatible with the acceleration of the Universe, with $\omega = -1$ and $\rho \propto 1$.
    \end{enumerate}
    The value of dust is justified by the fact that $p=0$ implies $\omega =0$. The value of radiation is justified by the fact that electromagnetic radiation is described by the traceless Maxwell tensor, which implies that the energy-momentum tensor must be traceless $\tr T = - \rho + 3 p = 0$ so that $p = \rho / 3$. Notice that they both agree with special relativistic considerations: for dust $E \sim m_0$ and $V \sim a^3$, whereas for radiation with red-shift $E \propto a^{-1}$. Furthermore, the dark energy contribution can be written in terms of the famous cosmological constant $\Lambda$
    \begin{equation*}
        \rho = - p = \frac{\Lambda}{8 \pi G_N} ~.
    \end{equation*}
    \begin{proof}
        In fact, using the Einstein field equations 
        \begin{equation*}
            G_{\mu\nu} + \Lambda g_{\mu\nu} = 8 \pi G_N T_{\mu\nu} ~, \quad G_{\mu\nu} = 8 \pi G_N T_{\mu\nu} - \Lambda g_{\mu\nu} ~,
        \end{equation*}
        \begin{equation*}
            G_{\mu\nu} = 8 \pi G_N ( T_{\mu\nu} - \frac{\Lambda g_{\mu\nu}}{8 \pi G_N} ) = 8 \pi G_N ( T_{\mu\nu} - T_{\mu\nu}^{\Lambda} ) ~,
        \end{equation*}
        hence, we have 
        \begin{equation*}
            T_{00}^{\Lambda} = \rho = \frac{\Lambda}{8 \pi G_N} = - p ~.
        \end{equation*}
    \end{proof}

\section{Friedmann equations}

    It is time to calculate the Einstein field equations for the FRW metric, which are called the Friedmann equations
    \begin{equation*}
        3 \Big ( \Big ( \frac{\dot a}{a} \Big)^2 + \frac{k}{a^2} \Big ) = 8 \pi G_N \rho ~, 
    \end{equation*}
    \begin{equation*}
        3 \frac{\ddot a}{a} = - 4 \pi G_N (\rho + 3 p) ~.
    \end{equation*}
    Notice that the curvature scalar $k$ appears only in the form $k / a^2$; in this way a rescaling is well-compensated. The first one is only a constraint to select initial conditions, while it is the second equation that governs the time evolution.
    \begin{proof}
        Because of isotropy, there are only $2$ independent equations. For the $00$-th components of the Einstein tensor
        \begin{equation*}
            G^0_{\phantom 0 0} = g^{00} G_{00} = - \frac{3}{a^2} (k + \dot a^2) ~,
        \end{equation*}
        we obtain 
        \begin{equation*}
            G^0_{\phantom 00} = 8 \pi G_N T^{0}_{\phantom 00} ~, \quad - \frac{3}{a^2} (k + \dot a^2) = - 8 \pi G_N \rho ~, 
        \end{equation*}
        \begin{equation*}
            3 \Big ( \Big ( \frac{\dot a}{a} \Big)^2 + \frac{k}{a^2} \Big ) = 8 \pi G_N \rho ~.
        \end{equation*}
        For the $11$-th components of the Einstein tensor
        \begin{equation*}
            G^1_{\phantom 11} = g^{11} G_{11} = - \frac{1}{a^2} (k + 2 a \ddot a + \dot a^2) ~,
        \end{equation*}
        we obtain 
        \begin{equation*}
            G^1_{\phantom 1} = 8 \pi G_N T^{1}_{\phantom 11} ~, \quad - \frac{1}{a^2} (k + 2 a \ddot a + \dot a^2) = 8 \pi G_N p ~.
        \end{equation*}
        Now, we use the first equation 
        \begin{equation*}
            \Big ( \frac{\dot a}{a} \Big)^2 + \frac{k}{a^2} = \frac{8}{3} \pi G_N \rho
        \end{equation*}
        to simplify the second one 
        \begin{equation*}
            - \frac{1}{a^2} (k + 2 a \ddot a + \dot a^2) = - \Big ( \frac{\dot a}{a} \Big)^2 - \frac{k}{a^2} - 2 \frac{\ddot a}{a} = - \frac{8}{3} \pi G_N \rho - 2 \frac{\ddot a}{a} ~.
        \end{equation*}
        Hence, we find
        \begin{equation*}
            - \frac{8}{3} \pi G_N \rho - 2 \frac{\ddot a}{a} = 8 \pi G_N p ~, \quad 6 \frac{\ddot a}{a} = - 24 \pi G_N p - 8 \pi G_N \rho ~,
        \end{equation*}
        \begin{equation*}
            3 \frac{\ddot a}{a} = - 4 \pi G_N (\rho + 3 p) ~.
        \end{equation*}
    \end{proof}
    \noindent Because of the Bianchi identity, the energy conservation law and the Friedmann equations are not independent.
    \begin{proof}
        In fact, deriving the first one
        \begin{equation*}
            \dv{}{t} \frac{\dot a^2}{a^2} + \dv{}{t} \frac{k}{a^2} = - 2 \frac{\dot a^2 }{a^3} + 2 \frac{\ddot a}{a^2} - 2 \frac{k}{a^3} = \frac{8}{3} \pi G_N \dot \rho ~,
        \end{equation*}
        then using the energy conservation law, we find 
        \begin{equation*}
            - 2 \frac{\dot a^2 }{a^3} + 2 \frac{\ddot a}{a^2} - 2 \frac{k}{a^3} = - 8 \pi G_N (\rho + p) \frac{\dot a}{a}  ~, \quad - \frac{\ddot a}{a} + \frac{\dot a^2 }{a^2}  + \frac{k}{a^2} = 4 \pi G_N (\rho + p) \dot a  ~, 
        \end{equation*}
        Finally, using the first one, we obtain 
        \begin{equation*}
            - \frac{\ddot a}{a} + \frac{8}{3} \pi G_N = 4 \pi G_N (\rho + p) \dot a  ~, \quad 3 \frac{\ddot a}{a} = - 4 \pi G_N (\rho + 3 p) ~.
        \end{equation*}
        Similarly, it can be shown that they are all related.
    \end{proof}

    It is useful to express the Friedmann equations in terms of the density parameter and critical density
    \begin{equation*}
        \Omega = \frac{8 \pi G_N \rho}{3 H^2} = \frac{\rho}{\rho_{crit}} ~, \quad \rho_{crit} = \frac{3 H^2}{8 \pi G_N} ~.
    \end{equation*}
    Therefore, the first Friedmann equation becomes 
    \begin{equation*}
        \Omega - 1 = \frac{k}{H^2 a^2} ~.
    \end{equation*}
    \begin{proof}
        Expressing the first equation in terms of $H$, we find
        \begin{equation*}
            \Big ( \frac{\dot a}{a} \Big)^2 + \frac{k}{a^2} = H^2 + \frac{k}{a^2} = \frac{8}{3} \pi G_N \rho ~, \frac{k}{a^2} = \frac{8}{3} \pi G_N \rho - H^2 ~,
        \end{equation*}
        \begin{equation*}
            \frac{k}{a^2 H^2} = \frac{8 \pi G_N \rho}{3 H^2} - 1 = \Omega - 1 ~.
        \end{equation*}
    \end{proof}
    \noindent As we said before, the sign of the curvature scalar $k$ can be used to determine which is the shape of the Universe. However, by means of this equation, we can equivalently use the density parameter $\Omega$ or the density $\rho$ in the following way
    \begin{enumerate}
        \item open universe if $k = -1$ or $\rho < \rho_{crit}$ or $\Omega < 1$, 
        \item flat universe if $k = 0$  or $\rho = \rho_{crit}$ or $\Omega = 1$, 
        \item closed universe if $k = 1$ or $\rho > \rho_{crit}$ or $\Omega > 1$.
    \end{enumerate}
    Now, we are able to solve the Friedmann equations for the simple flat case: 
    \begin{enumerate}
        \item flat matter-dominated, for which
            \begin{equation*}
                k = 0 ~, \quad \rho \sim a^{-3} \quad \Rightarrow \quad a \sim t^{2/3} ~;
            \end{equation*}
        \item flat radiation-dominated, for which
            \begin{equation*}
                k = 0 ~, \quad \rho \sim a^{-4} \quad \Rightarrow \quad a \sim t^{1/2} ~;
            \end{equation*}
        \item flat dark energy-dominated, for which
            \begin{equation*}
                k = 0 ~, \quad \rho \sim \Lambda \quad \Rightarrow \quad a \sim \exp(H_0 t) ~.
            \end{equation*}
    \end{enumerate}

    \begin{proof}
        For the flat matter-dominated, i.e.~$k = 0$ or $\Omega = 1$ and $\rho \sim a^{-3}$, we have 
        \begin{equation*}
            1 = \Omega \sim \frac{\rho}{H^2} = \frac{1}{a^3} \frac{a^2}{\dot a^2}= \frac{1}{a \dot a^2} ~, \quad a \dot a^2 = 1 ~,
        \end{equation*}
        \begin{equation*}
            a^{1/2} \dot a \sim 1 ~, \quad a^{1/2} da \sim dt ~, \quad a^{3/2} \sim t ~, \quad a \sim t^{2/3} ~.
        \end{equation*}
        For the flat radiation-dominated, i.e.~$k = 0$ or $\Omega = 1$ and $\rho \sim a^{-4}$, we have 
        \begin{equation*}
            1 = \Omega \sim \frac{\rho}{H^2} = \frac{1}{a^4} \frac{a^2}{\dot a^2}= \frac{1}{a^2 \dot a^2} ~, \quad a^2 \dot a^2 = 1 ~,
        \end{equation*}
        \begin{equation*}
            a \dot a \sim 1 ~, \quad a da \sim dt ~, \quad a^2 \sim t ~, \quad a \sim t^{1/2} ~.
        \end{equation*}
        For the flat dark energy-dominated, i.e.~$k = 0$ or $\Omega = 1$ and $\rho \sim \Lambda$, we have 
        \begin{equation*}
            1 = \Omega \sim \frac{\rho}{3 H^2} = \frac{\Lambda}{3} \frac{a^2}{\dot a^2} ~, \quad \dot a^2 = \frac{\Lambda}{3} a^2 ~,
        \end{equation*}
        \begin{equation*}
            \dot a \sim (\Lambda/3)^{1/2} a ~, \quad \frac{da}{a} \sim (\Lambda/3)^{1/2} dt ~, \quad a \sim \exp ( (\Lambda/3)^{1/2} t ) = \exp ( H_0 t ) ~.
        \end{equation*}
    \end{proof}

    We conclude the chapter with a comparison between the Einstein and the Newton cosmology. Consider a particle of mass $m$ on the surface of radius $R = r a$. The Newtonian second law reads as 
    \begin{equation*}
        \ddot R = r \ddot a = -V_N = - \frac{G_N M}{R^2} = - \frac{G_N \rho}{R^2} V = - \frac{G_N \rho}{R^2} \frac{4}{3} \pi R^3 = - \frac{4}{3} \pi G_N \rho R = - \frac{4}{3} \pi G_N \rho r a ~, 
    \end{equation*}
    \begin{equation*}
        3 \frac{\ddot a}{a} = - 4 \pi G_N \rho ~,
    \end{equation*}
    which shows that the Newtonian limit can be reached only if $p = 0$.

\chapter{Expansion of the universe}

    Hubble's law states that galaxies are moving away from us at a velocity that is proportional to their distance. It is important because it gives a direct observation of the Hubble constant. In this chapter, we will study how to derive this law from observational quantities: luminosity and red-shift. In the final section, we will briefly review what is called the standard model of cosmology: cosmic background radiation, inflation and $\Lambda$-CDM model.

\section{Cosmological red-shift}

    Place the origin at $r=0$. Consider a comoving source at $r_1$ and a comoving receiver at $r_2$. Two emitted radially ($\theta$ and $\phi$ constant) electromagnetic waves at times $t_1$ and $t_1 + \delta t_1$ will be received at times $t_2$ and $t_2 + \delta t_2$. We want to calculate the difference in frequency between them. Since light has zero interval, we have 
    \begin{equation*}
        0 = ds^2 = - dt^2 + \frac{a^2(t)}{1 - kr^2} dr^2 ~, \quad dt = \frac{a(t)}{\sqrt{1 - kr^2}} dr ~.
    \end{equation*}
    For the first signal, we have 
    \begin{equation*}
        \int_{r_1}^{r_2} \frac{dr}{\sqrt{1 - k r^2}} = \int_{t_1}^{t_2} \frac{d t}{a(t)} ~,
    \end{equation*}
    whereas, for the second signal, we have
    \begin{equation*}
        \int_{r_1}^{r_2} \frac{dr}{\sqrt{1 - k r^2}} = \int_{t_1 + \delta t_1}^{t_2 + \delta t_2} \frac{dt}{a(t)} ~.
    \end{equation*}
    Equating them two, we find 
    \begin{equation*}
        \int_{t_1}^{t_2} \frac{d t}{a(t)} = \int_{t_1 + \delta t_1}^{t_2 + \delta t_2} \frac{dt}{a(t)} ~,
    \end{equation*}
    \begin{equation*}
        \int_{t_1}^{t_1 + \delta t_1} \frac{d t}{a(t)} + \cancel{\int^{t_2}_{t_1 + \delta t_1} \frac{d t}{a(t)}} = \cancel{\int_{t_1 + \delta t_1}^{t_2} \frac{dt}{a(t)}} + \int_{t_2}^{t_2 + \delta t_2} \frac{dt}{a(t)} ~,
    \end{equation*}
    \begin{equation*}
        \int_{t_1}^{t_1 + \delta t_1} \frac{d t}{a(t)} = \int_{t_2}^{t_2 + \delta t_2} \frac{dt}{a(t)} ~.
    \end{equation*}
    Making the reasonable assumption that $\delta t_1$ and $\delta t_2$ are so small that the cosmic scale factor $a$ does not change significantly, then we obtain
    \begin{equation*}
        \frac{\delta t_1}{a(t_1)} = \frac{\delta t_2}{a(t_2)} ~, \quad \frac{\delta t_1}{\delta t_2} = \frac{a(t_1)}{a(t_2)} ~.
    \end{equation*}
    It is useful to rewrite in terms of frequencies
    \begin{equation*}
        \frac{\omega(t_1)}{\omega(t_2)} = \frac{2\pi}{\delta t_1} \frac{\delta t_2}{2\pi} = \frac{\delta t_2}{\delta t_1} = \frac{a(t_2)}{a(t_1)} ~,
    \end{equation*}
    or in terms of the red-shift $z$, defined as
    \begin{equation*}
        z = \frac{\lambda_2 - \lambda_1}{\lambda_1} = \frac{a(t_2)}{a(t_1)} - 1 ~,
    \end{equation*}
    where $\lambda_2$ is the observed wavelength and $\lambda_1$ is emitted observed. This red-shift is from different nature than Doppler effect, since it is caused by the expansion of spacetime. If the Universe is expanding, i.e.~$a(t_2) > a(t_1)$, the red-shift is positive, otherwise if the Universe is contracting, i.e.~$a(t_2) < a(t_1)$, the red-shift is negative, which is called blue-shift. It is symmetric between observer and receiver, unlike the gravitational red-shift. Furthermore, it is important highlight that $a(t)$ can be rescaled, therefore is not physical, but in our formula it is the ratio that compares. It depends only on time of emission and arrival (it is not cumulative); it is due a different reference frame used by them two.

\section{Luminosity-distance relation}

    In astronomy, the only physical quantity we can directly measure is the apparent luminosity of a stellar object. However, we can measure also distance using the luminosity-distance relation.
    
    Consider first the Minkowski metric. Since energy is conserved during propagation, the flux of energy through a surface of radius $R$ is time-independent and depends only on radius 
    \begin{equation*}
        F = \frac{L}{A} = \frac{L_0}{4 \pi R^2} ~.
    \end{equation*}
    Therefore, the luminosity-distance relation is 
    \begin{equation*}
        d_L^2 = R^2 = \frac{L_0}{4 \pi F} ~. 
    \end{equation*}

    Consider now the FRW metric. In the previous section, we have seen that photons are red-shifted by a factor $(1 + z)$. In particular, the observer receives a signal that is delayed by a factor $(1+z)^2$ and the luminosity transforms as
    \begin{equation*}
        L \simeq \frac{L_0}{(1 + z)^2} ~, 
    \end{equation*} 
    hence, the flux becomes
    \begin{equation*}
        F = \frac{L}{A} = \frac{L}{4 \pi R^2} = \frac{L}{4 \pi(a_0 r)^2} \simeq \frac{L_0}{4 \pi(a_0 r)^2 (1+z)^2} = \frac{L_0}{4 \pi d_L^2} ~,
    \end{equation*} 
    where the luminosity-distance relation is given by
    \begin{equation}
        d_L = a_0 r (1 + z) \simeq \sqrt{\frac{L_0}{F}} ~.
    \end{equation}
    Notice that for $z > 0$, we have darker objects. Besides, in order to measure the red-shift, we need the emitted original frequency. In astronomy, this can be achieved by using standard candles: characteristic behaviours, like period of oscillation of luminosity. 

\section{Hubble's law}

    After we have introduced these two observable quantities, we are able to extract the Hubble law. The first assumption is that we are working with small red-shift so that we can work in power series in $z$. Therefore, we expand the cosmic scale factor in time of emission $t=t_1$ around $t_0$
    \begin{equation*}
        a(t) = a_0 + \dot a_0 (t-t_0) + \frac{\ddot a_0}{2} (t-t_0)^2 + O((t-t_0)^3) ~.
    \end{equation*}
    We introduce the deceleration parameter 
    \begin{equation*}
        q = - \frac{a \ddot a}{\dot a^2} 
    \end{equation*}
    and rewrite the expansion in terms of $q_0$ and $H_0$
    \begin{equation*}
    \begin{aligned}
        a(t) & = a_0 \Big ( 1 + \frac{\dot a_0}{a_0} (t-t_0) + \frac{1}{2} \frac{\dot a_0^2}{a_0^2} \frac{\ddot a_0 a_0}{\dot a_0^2} (t-t_0)^2 \Big ) + O((t-t_0)^3) \\ & = a_0 \Big ( 1 + H_0 (t-t_0) - \frac{1}{2} H_0^2 q_0 (t-t_0)^2 \Big ) + O((t-t_0)^3) ~.
    \end{aligned}
    \end{equation*}
    Therefore, the red-shift can be expanded as 
    \begin{equation*}
    \begin{aligned}
        z & = \frac{a_0}{a} - 1 \simeq \Big ( 1 + H_0 (t-t_0) - \frac{1}{2} H_0^2 q_0 (t-t_0)^2 \Big)^{-1} - 1 + O((t-t_0)^3) \\ & \simeq 1 - H_0 (t-t_0) + \frac{1}{2} H_0^2 q_0 (t-t_0)^2 + H_0^2 (t-t_0)^2 - 1 + O((t-t_0)^3) \\ & = - H_0 (t-t_0) + H_0^2 \Big ( \frac{q_0}{2} + 1 \Big ) (t-t_0)^2 + O((t-t_0)^3) ~,
    \end{aligned}
    \end{equation*}
    where we have used the expansion 
    \begin{equation*}
        \frac{1}{1 + x} = 1 - x + x^2 + O(x^3) ~.
    \end{equation*}
    Now, we need to invert this relation in order to find $t_0-t$ in terms of $z$. Through an ansatz 
    \begin{equation*}
        t_0 - t = A z + B z^2 + O(z^3) ~,
    \end{equation*}
    we find 
    \begin{equation*}
        t_0 - t = A H_0 (t_0 - t) + A H_0^2 \Big ( \frac{q_0}{2} + 1 \Big ) (t_0 - t)^2 + B H_0^2 (t_0 - t)^2 ~,
    \end{equation*}
    which implies 
    \begin{equation*}
        A H_0 = 1 ~, \quad A = \frac{1}{H_0} ~,
    \end{equation*}
    \begin{equation*}
        A H_0^2 \Big ( \frac{q_0}{2} + 1 \Big ) + B H_0^2 = 0 ~, \quad \Big ( \frac{q_0}{2} + 1 \Big ) + B H_0  = 0 ~,
    \end{equation*}
    \begin{equation*}
        B H_0 = - \frac{q_0}{2} - 1 ~, \quad B = - \frac{1}{H_0} \Big (\frac{q_0}{2} + 1 \Big ) ~.
    \end{equation*}
    Putting together, we obtain 
    \begin{equation*}
        t_0 - t = \frac{1}{H_0} \Big ( z - \Big (\frac{q_0}{2} + 1 \Big ) z^2 \Big ) + O(z^3) ~.
    \end{equation*}

    Our final goal is to find the Hubble's law, which is written in terms of distance
    \begin{equation*}
    \begin{aligned}
        \int_{t}^{t_0} \frac{dt}{a(t)} = \int_0^r \frac{dr}{\sqrt{1 - kr^2}} \simeq \int_0^r dr = r ~.
    \end{aligned}
    \end{equation*}
    Therefore, we have
    \begin{equation*}
    \begin{aligned}
        r & \simeq \int_{t}^{t_0} \frac{dt}{a(t)} \simeq \int_{t}^{t_0} dt ~ \frac{1}{a_0} (1 + H_0 (t - t_0))^{-1} \simeq \int_{t}^{t_0} dt ~ \frac{1}{a_0} (1 - H_0 (t - t_0)) \\ & = \int_{t}^{t_0} dt ~ \frac{1}{a_0} (1 + H_0 t_0 - H_0 t) = \frac{1}{a_0} (t + H_0 t_0 t - H_0 \frac{t^2}{2}) \Big \vert_{t}^{t_0} \\ & = \frac{1}{a_0} (t_0 + H_0 t_0^2 - H_0 \frac{t_0^2}{2} - t - H_0 t_0 t + H_0 \frac{t^2}{2} ) + O((t-t_0)^3) \\ & = \frac{1}{a_0} \Big ( (t_0 - t) + H_0 \frac{t_0^2}{2} - H_0 t_0 t + H_0 \frac{t^2}{2}\Big ) + O((t-t_0)^3) \\ & = \frac{1}{a_0} \Big ( (t_0 - t) + \frac{H_0}{2} (t_0 - t)^2 \Big ) + O((t-t_0)^3) ~.
    \end{aligned}
    \end{equation*}
    Finally, combining the two expressions, we obtain 
    \begin{equation*}
        r \simeq \frac{1}{H_0 a_0} \Big (  z - (\frac{q_0}{2} + 1  ) z^2 + \frac{1}{2} z^2 \Big ) + O(z^3) = \frac{1}{H_0 a_0} \Big (  z + \frac{1}{2} (1 - q_0) z^2 \Big ) + O(z^3)
    \end{equation*}
    and the distance-luminosity relation becomes 
    \begin{equation*}
    \begin{aligned}
        d_L & = a_0 r (1+z) = a_0 (1+z) \frac{1}{H_0a_0} \Big (z + \frac{1}{2} (1 - q_0) z^2 \Big ) + O(z^3) \\ & \simeq \frac{1}{H_0} \Big (z + \frac{1}{2} (1 - q_0) z^2 \Big ) + O(z^3) ~,
    \end{aligned}
    \end{equation*}
    which is indeed the so-called Hubble's law. Physically, it means that galaxies recede from us at a velocity $v \sim z$ proportional to the distance $d_L$.

\section{Standard model of cosmology}

    The solutions of the Friedmann equations shows that the Universe must have started from a singularity at $a = 0$ and expands from that ever since. This can be seen by writing the second Friedmann equation in terms of the deceleration parameter
    \begin{equation*}
        q = \frac{4 \pi G_N}{3 H^2} (\rho + 3 p) ~.
    \end{equation*}
    For vanishing cosmological constant, we have $q > 0$ or $\ddot a < 0$, hence the universe is decelerating its expansion due to gravitational attraction. Furthermore, given that $a > 0$ (by definition), $\dot a > 0$ (red-shift) and $\ddot a < 0$, $a$ does not have a turning point and it must have reached $a = 0$ in finite time. For a radiation-dominated universe, the energy density goes like $\rho \sim a^{-4}$, which menas that for $a \rightarrow 0$, it goes to infinity. This is the basis of the Hot Big Bang model. Its main adversary was the steady model, in which energy density remains constant. However, observations on the Cosmic Microwave Background radiation shows that the Universe is cooling down. The CMB radiation is the almost relic radiation generated by the decoupled of radiation from matter on the last scattering surface. Before, photons were in thermal equilibrium with electrons and protons, and the only way to understand the physics of that period is through gravitational waves. Another important observational feedbacks from the CMB are the fact that its fine anisotropies shows that the Universe is flat, but there is an high degree of isotropy. The latter is very suspiciuos, because it means that the light we receive from different direction in the sky should not have the same temperature, unless they were causally connected in the past. How to solve this problem, called the Horizon problem?

    Suppose we place ourselves at $r = 0$ and $t=0$. We want to find the proper distance that the photon has travelled from $r = r_0$ and $t= - t_0$ to us
    \begin{equation*}
        R \sim a( - t_0) r_s ~.
    \end{equation*}
    Using the expression above, we need to evaluate the integral
    \begin{equation*}
        r_0 = \int_0^{r_0} dr \sim \int_{-t_0}^0 \frac{dt}{a(t)} ~.
    \end{equation*}
    Consider first a matter/radiation-dominated Universe, so that
    \begin{equation*}
        a(t) \sim t^\alpha ~, \quad r_0 \sim \int_{-t_0}^0 dt ~ t^{- \alpha} \sim t_0^{1 - \alpha} ~,
    \end{equation*}
    where $\alpha = 2 / 3, 1 / 2$. Therefore, the proper distance is given by
    \begin{equation*}
        R \sim a(-t_0) r_0 = t^\alpha_0 t^{1-\alpha}_0 = t_0 ~,
    \end{equation*}
    which coincides with Minkowski. Furthermore, 
    \begin{equation*}
        \dot a \sim t^{\alpha - 1} ~, \quad H = \frac{\dot a}{a} \sim t^{-1} ~,
    \end{equation*}
    and the particle horizon grows linearly with time  
    \begin{equation}
        R_H \sim \frac{1}{H(-t_0)} \sim t_0 ~.
    \end{equation}
    However, in a dark energy-dominated Universe, we have
    \begin{equation*}
        a(t) \sim \exp(H_0 t) ~, \quad r_s \sim \frac{\exp(-H_0 t)}{H_0} ~,
    \end{equation*}
    and the particle horizon is constant
    \begin{equation*}
        R \sim \exp(H_0 t_0) r_s \sim \frac{1}{H_0} ~.
    \end{equation*}
    This can explain CMB homogeneity: the Universe started so small to have all particles in causal connection and, then, there was an inflation phase of rapid expansion, that eventually ended.

    We conclude this chapter with experimental results. As said before, the CMB provides us that the universe is flat with $\Omega \simeq 1$ and density $\rho_0 = \rho_{crit} \simeq 10^{-29} g/cm^3$. The contribution to density can be identified by the following sources 
    \begin{enumerate}
        \item radiation is negligible with $\rho_{rad} \simeq 0$
        \item massive baryonic matter with $\rho_{mat} / \rho_0 \simeq 5 \%$
        \item dark matter with $\rho_{DM} / \rho_0 \simeq 25 \%$
        \item dark energy with $\rho_{DE} / \rho_0 \simeq 70 \%$
    \end{enumerate}
    Dark matter is needed to explain the rotation of galaxies and dark energy to explain the expansion of the Universe.

